% MAIN TEXT, v8 Section 4. Stage 4's infinite-horizon benchmark will follow.
\section{Persistence: growing versus bounded budgets}
\label{sec:persistence}

\subsection{Unbounded physical budgets}

\begin{proposition}[Growth obstruction and sharp productivity scaling]
\label{prop:growth-trichotomy}
Let $B(t)\to\infty$, let $\ell$ be concave, and retain the assumptions
of \cref{prop:hub}. Suppose that eventually some direction has
$p_j(t)\leq M<\infty$ and $S_j(t)\geq s>0$.
For $\eta\geq1$, the all-inward corner is excluded at every sufficiently
late time, for every path of finite positive productivities $a(t)$.
For $0\leq\eta<1$, eventual inwardness instead requires
\begin{equation}
 \Omega(t)\geq
 \left[\frac{s\Gamma(B(t))}{M}\right]^{1/(1-\eta)}
 \geq\left[\frac{sB(t)}{M}\right]^{1/(1-\eta)}.
 \label{eq:necessary-output-growth}
\end{equation}
For the sufficiency statement, additionally fix all $p_j>0$, finite
$r_j,S_j\geq0$, and $d_0$. If $0<\eta<1$, it is sufficient that,
eventually,
\begin{equation}
 \Omega(t)\geq c_\Omega\Gamma(B(t))^{1/(1-\eta)},\qquad
 c_\Omega>
 \left(\max_j\frac{S_j}{p_j}\right)^{1/(1-\eta)}.
 \label{eq:sufficient-output-growth}
\end{equation}
The inequalities then become strict. For $\eta=0$, the exact condition is
\begin{equation}
 \frac{\Omega(t)}{\Gamma(B(t))}=a(t)\ell'(B(t))
 \geq\max_j\frac{r_j+S_j}{p_j}.
 \label{eq:linear-value-growth}
\end{equation}
In particular, for $\ell(z)=z$ and $0<\eta<1$,
$a(t)\geq c_\Omega B(t)^{\eta/(1-\eta)}$ with the strict coefficient
condition above is eventually sufficient.
\end{proposition}
\noindent\emph{Proof sketch.}
For $\eta=1$ the residual term has positive infimum one after division
by $S_j\Gamma$. For $\eta>1$ that infimum is
\begin{equation}
 \min_{\Omega>0}\frac{(d_0+\Omega)^\eta}{\Omega}
 =d_0^{\eta-1}\frac{\eta^\eta}{(\eta-1)^{\eta-1}}.
 \label{eq:residual-minimum}
\end{equation}
Both conflict with $\Gamma(B)\geq B\to\infty$. Below unit curvature,
comparison with $\Omega^{\eta-1}$ gives the necessary rate; the strict
coefficient in \cref{eq:sufficient-output-growth} leaves an asymptotic
margin for remote computational returns. See \cref{app:growth-proof}.

This result concerns the growth of physical resources, not the growth
of computational efficiency. With persistent residual value and bounded
prices, productivity can compensate for an unbounded resource budget
only on the $\eta<1$ branch of the concave-production model. The endpoint
$\eta=0$ must retain $r_j$ as well as $S_j$; remote computation is not a
vanishing correction in that case.

\begin{corollary}[Logarithmic pathwise specialization]
\label{prop:persistence}
Let $\ell(z)=\ln(1+z)$ and $\eta=1$, with fixed $r_j,d_0$ and finite
positive $B(t),a(t),p_j(t)$. Define $H_j(t),A_j(t)$ by
\cref{eq:HA}. On any specified path, the corner is optimal at every
time exactly when, in every direction with $A_j(t)>0$,
\begin{equation}
 p_j(t)>H_j(t),\qquad
 a(t)\geq\max_{j:A_j(t)>0}
       \frac{A_j(t)}{p_j(t)-H_j(t)}.
 \label{eq:pathwise}
\end{equation}
In particular, where $S_j(t)>0$,
\begin{equation}
 \frac{p_j(t)}{S_j(t)}>[1+B(t)]\ln[1+B(t)].
 \label{eq:path-ratio}
\end{equation}
For $B(t)\to\infty$, fixed positive residual slopes therefore require
prices at least of order $B\ln B$; bounded prices require
$S_j(t)=O((B\ln B)^{-1})$. These rate restrictions do not replace the
productivity inequality. A zero-slope direction with $r_j>0$ requires
$a(t)\geq[1+B(t)]r_j/p_j(t)$.
\end{corollary}
\begin{proof}
Apply \cref{eq:exact,eq:acrit} at each time. Non-negativity and $d_0>0$
imply $A_j=0$ exactly when $r_j=S_j=0$. Division by a positive slope
and $(1+B)\ln(1+B)\sim B\ln B$ give the rate statements.
\end{proof}

% MOVE HERE or immediately after this subsection, unchanged:
% the pure-computation paragraph containing eq:Rstar and eq:budgetcrit,
% and Figure 2 (fig:budget). Its algebraic reduction is also recorded
% in the appendix proof of prop:growth-trichotomy.

\subsection{What increasing returns would have to accomplish}

\begin{proposition}[At least exponential local returns]
\label{prop:exponential-necessity}
Let $\eta=1$ and retain the model assumptions apart from local
concavity. Let a fixed local production function satisfy $\ell(0)=0$
and be continuously differentiable and positive with positive derivative
for positive budgets. Let $B(t)$ be continuous with $B(t)\to\infty$.
If eventually $p_j(t)\leq M$, $S_j(t)\geq s>0$, and the corner is
optimal at every time, then for all sufficiently large $B$,
\begin{equation}
 \frac{\dd\ln\ell(B)}{\dd B}>\frac{s}{M},\qquad
 \ell(B)>\ell(B_0)\exp\!\left[\frac{s}{M}(B-B_0)\right]
 \quad(B>B_0).
 \label{eq:exponential-necessity}
\end{equation}
Convex local production gives $\Gamma(B)\leq B$, but polynomial
superlinearity is insufficient: $\ell(z)=z^k$ gives $\Gamma=B/k$.
By contrast, $\ell(z)=[\exp(\psi z)-1]/\psi$, $\psi>0$, gives
$\Gamma=[1-\exp(-\psi B)]/\psi<1/\psi$.
\end{proposition}
\noindent\emph{Proof sketch.}
The necessary strict floor gives $\ell'/\ell>S_j/p_j\geq s/M$.
Continuity of the budget path covers every sufficiently large budget,
so integration proves the exponential bound. The remaining identities
follow by differentiation; see \cref{app:returns-proof}.

The required increase is with local resource investment, not merely
with calendar time. This supplies a quantitative test for
compression-based inward development proposals such as
\citet{Smart2012}. Increasing returns remove the concavity guarantee.
For exponential local production with one affine external channel and
linear residual value, an exact global test is available: the corner
must satisfy its first-order condition and outperform the all-external
endpoint (\cref{lem:exponential-endpoints}). Smooth concave residual
utilities, or smooth concave external production, can invalidate this
two-test shortcut. Explicit counterexamples are given in
\cref{app:exponential-counterexamples}.

\subsection{The tipping budget}

\begin{proposition}[A unique productivity-independent tipping budget]
\label{prop:tipping}
Let $\eta=1$ and $\ell\in C^2([0,\infty))$, with $\ell(0)=0$,
$\ell'>0$, and $\ell''\leq0$. Fix $d_0\geq1$, all $p_j>0$, and finite
$r_j,S_j\geq0$. If at least one $S_j>0$, define $B_{\rm tip}$ by
\begin{equation}
 \Gamma(B_{\rm tip})=\min_{j:S_j>0}\frac{p_j}{S_j}.
 \label{eq:tipping-budget}
\end{equation}
This budget exists and is unique, since
\begin{equation}
 \Gamma(0)=0,\qquad
 \Gamma'(B)=1-\frac{\ell(B)\ell''(B)}{\ell'(B)^2}\geq1.
 \label{eq:gamma-monotone}
\end{equation}
The corner is attainable at finite productivity exactly for
$0<B<B_{\rm tip}$, where its threshold is
$a_{\rm crit}^{\ell}(B)$ from \cref{eq:acrit-general}.
This threshold is non-decreasing and diverges as
$B\uparrow B_{\rm tip}$. At $B=B_{\rm tip}$ as well as above it,
no finite productivity makes the corner optimal.
If eventually $B(t)\leq B'<B_{\rm tip}$ and $a(t)\to\infty$, every
static optimizer is eventually all-inward. If every $S_j=0$, set
$B_{\rm tip}=\infty$; no finite productivity-independent residual
barrier exists.
\end{proposition}
\noindent\emph{Proof sketch.}
Differentiate $\Gamma$. Below its unique crossing, every denominator
in \cref{eq:acrit-general} is positive. Each nontrivial threshold is
non-decreasing because its denominator has derivative
$p_j\ell''-S_j\ell'\leq0$. A binding positive-slope denominator
vanishes at the crossing. See \cref{app:tipping-proof}.

\begin{corollary}[Fixed- and bounded-budget curvature]
\label{prop:curvature}
At a fixed positive budget and under concave production, each nontrivial
direction has a unique finite high-productivity threshold for
$0\leq\eta<1$. At $\eta=1$, \cref{eq:acrit-general} applies.
At $\eta>1$, a direction with $S_j>0$ has an inward set in $a>0$
that is empty, a singleton, or a bounded closed interval. A direction
with $S_j=0$ requires only $a\geq r_j/[p_j\ell'(B)]$ for every $\eta$.
For the local logarithm the test is precisely
\begin{equation}
 r_j+S_j(d_0+aL)^\eta\leq\frac{p_ja}{1+B},
 \qquad L=\ln(1+B).
 \label{eq:iso-corner}
\end{equation}
For the path statements assume additionally that $\ell$ is $C^1$ with
positive derivative near the positive limiting budget. Consequently, if
$B(t)\to\bar B\in(0,\infty)$ and $a(t)\to\infty$,
fixed parameters give eventual inwardness for $\eta<1$. More generally, this conclusion holds for any $0<B(t)\leq B_{\max}<\infty$ if $\ell$ is concave and $\ell'(B_{\max})>0$. For $\eta>1$,
a positive residual slope excludes eventual inwardness. At $\eta=1$,
$\bar B<B_{\rm tip}$ gives eventual inwardness and
$\bar B>B_{\rm tip}$ excludes it. At equality, the rate of approach
from below must be compared with $a_{\rm crit}^{\ell}(B(t))$.
\end{corollary}
\noindent\emph{Proof.}
See \cref{app:curvature-proof}. This is the fixed-budget curvature
classification, with a uniform bound near the positive limiting budget;
it is not a new assumption about preferences. ``No tipping budget''
for $\eta<1$ means no productivity-independent finite obstruction,
not absence of a budget threshold at a fixed productivity.

\subsection{Resource acquisition and the inward basin}

Separate resources available locally from those supplied by prior
replicating acquisition. In discrete planning periods, write
\begin{equation}
 B_t=B_{\rm loc}(t)+g_XX_t,\qquad
 X_{t+1}=X_t+x_t,\qquad g_X\geq0,
 \label{eq:acquisition-stock}
\end{equation}
where $B_{\rm loc}(t)$ is non-decreasing with limit $\bar B>0$.
Only the replicating channel enters $X$; observation alone does not.
With no previous acquisition, an all-inward history has $X=0$ and
$B_t\leq\bar B$. Stopping new investment later does not erase an
inherited external stock. Since both stock and local mobilization are
non-decreasing, a crossing of $B_{\rm tip}$ cannot be reversed in this
model. In a multi-channel model this excludes all-inward expenditure,
not necessarily zero replication.

A basin cannot be specified by four scalar parameters without also
specifying the resource-mobilization path and the full residual utility.
The following benchmark fixes both. It uses one replicating channel,
linear residual value, and an already mobilized local ceiling.

\begin{proposition}[A unique inward-basin boundary]
\label{prop:basin}
Fix $p>0$, $S>0$, $r\geq0$, $d_0\geq1$, $g_a>0$, and $g_X>0$.
Let $X_0=0$, $B_{\rm loc}(t)=\bar B$, $a_t=a_0\exp(g_at)$, and
$0<\bar B<B_{\rm tip}$, where
$(1+B_{\rm tip})\ln(1+B_{\rm tip})=p/S$.
At each period choose the unique maximizer
\begin{equation}
 x^*(B,a)=\underset{0\leq x\leq B/p}{\arg\max}
 \left\{\ln[d_0+a\ln(1+B-px)+rx]+Sx\right\}.
 \label{eq:basin-policy}
\end{equation}
There is a unique
$a_{\rm sep}(g_a,g_X,\bar B)\in(0,a_{\rm crit}^{\ell}(\bar B))$
with the following classification:
\begin{itemize}[leftmargin=*,nosep]
\item If $a_0>a_{\rm sep}$, new external investment stops in finite time
and stays zero, at a budget strictly below $B_{\rm tip}$.
\item If $a_0<a_{\rm sep}$, the budget reaches or exceeds $B_{\rm tip}$
in finite time and new external investment remains positive thereafter.
\item If $a_0=a_{\rm sep}$, every finite-time investment is positive,
but $B_t\uparrow B_{\rm tip}$ and
$g_X\sum_{t\geq0}x_t=B_{\rm tip}-\bar B$.
\end{itemize}
Whenever $B_t\geq B_{\rm tip}$,
\begin{equation}
 x_t>\frac{B_t-B_{\rm tip}}p,\qquad
 B_{t+1}-B_{\rm tip}>
 (1+g_X/p)(B_t-B_{\rm tip}).
 \label{eq:runaway-bound}
\end{equation}
The boundary $a_{\rm sep}$ is non-increasing in $g_a$ and
non-decreasing in $g_X$ and in $\bar B/B_{\rm tip}$, with the other
parameters fixed.
\end{proposition}
\noindent\emph{Proof sketch.}
The unique static investment is continuous, increases with the budget,
and decreases with productivity. A path that neither stops nor crosses
must converge to the tipping budget. Two such paths cannot coexist:
a positive gap between their budgets cannot shrink. Continuity separates
the open sets of stopping and crossing paths, giving exactly one boundary.
The complete proof and finite-time basin certificates are in
\cref{app:basin-proof}.

For example, let $(p,S,r,d_0)=(1,0.002,0.05,1)$,
$g_a=0.05$, $g_X=0.5$, and $\bar B=0.6B_{\rm tip}$. New calculations
give $B_{\rm tip}=106.001289$,
$a_{\rm crit}^{\ell}(\bar B)=7.279612$, and
$a_{\rm sep}\simeq6.63216$. Starting at $a_0=0.5$ crosses the tipping
budget at $t=2$; starting at $a_0=50$ is inward from the start. With
$a_0=0.5$ but $g_X=0$, investment stops at $t=54$, not at $t=0$.
Thus removing resource feedback removes this lock-in mechanism, not
the initial incentive to expand.

Lock-in here names an absorbing allocation regime, not a welfare
mistake. Additional resources do not reduce the feasible set. In this
costless, non-depreciating stock benchmark, any finite-horizon
continuation value is non-decreasing in $X$: this follows by backward
induction because a larger stock permits the same current project and
more local investment. A strictly improving myopic external deviation
therefore remains improving after adding such a continuation. A
forward-looking planner need not avoid the tipping budget. Productive
local alternatives require the separate dynamic comparison developed
in the infinite-horizon subsection.

\paragraph{Speculative physical mapping (conjecture).}
If a fixed-horizon energy budget increases with usable host luminosity,
and if denser stellar neighbourhoods reduce external resource prices
while leaving residual values fixed, both changes raise
$\bar B/B_{\rm tip}$. The basin comparison then suggests a greater
propensity to expand. The conjecture concerns these physical parameter
mappings, not the proved comparative static.
